\documentclass[10pt,a4paper]{amsart}
\usepackage[margin=19mm]{geometry}
\usepackage{amsmath,amssymb,mathtools}
\usepackage[scr=rsfs,cal=euler]{mathalfa}
\usepackage{xfrac}
\usepackage[unicode,hidelinks,bookmarks=false]{hyperref}
\newcommand{\ie}{i.e.\ }
\newcommand{\cf}{cf.\ }
\newcommand{\resp}{respectively\ }
\newcommand{\Subsection}{\subsection}
\providecommand{\nobreakdash}{\nobreak\hbox{-}}
\setlength{\emergencystretch}{2em}
\multlinegap0pt
\usepackage{amssymb}
\usepackage{mathtools}
\usepackage[shortlabels]{enumitem}
\usepackage{xcolor}
\newtheorem{theorem}{Theorem}
\newtheorem{corollary}{Corollary}
\newtheorem{lemma}{Lemma}
\newcommand{\MW}{\mathrm{MW}}
\DeclareMathOperator{\Ric}{Ric}
\newcommand{\Sph}{\mathbb S}
\newcommand{\ent}{\mathrm{ent}}
\newcommand{\eps}{\varepsilon}
\title{On the proof of Bray's conjecture}
\author{Xumin Jiang, Mingxiang Li, Zhehui Wang}
\date{Source: arXiv:2608.20215v2 (CC BY 4.0)}
\begin{document}
\maketitle

\noindent\emph{Source and licence.} On the proof of Bray's conjecture. Version arXiv:2608.20215v2 (CC BY 4.0), distributed by arXiv under the Creative Commons Attribution 4.0 International licence, http://creativecommons.org/licenses/by/4.0/, as recorded in the arXiv licence metadata for this exact version and shown on the abstract page https://arxiv.org/abs/2608.20215v2. Full licence notice: \texttt{licenses/arxiv-2608.20215-CC-BY-4.0.txt}.\par\smallskip

\noindent\textit{Editorial scope.} Counted unit: the article's theorem theorem (source0.tex lines 681--683). The article's complete original proof of it is delivered with it (source0.tex lines 684--753, one \texttt{proof} environment of 70 lines). No mathematics is written, repaired or re-derived: the statement and the proof are the article's own lines, in the article's own notation, and no retained line is substituted or re-flowed.  Retained uncounted as the source-local context the proof needs: lines 99--101; lines 107--109; lines 190--194; lines 319--336; lines 322--324; lines 468--472; lines 524--539; lines 547--553; lines 578--580.  Nothing was omitted from any retained span, so the package carries no omission record.  No retained line contains a citation, so the wrapper carries no bibliography and no bibliography entry.  No retained line contains a figure environment, an image-inclusion command or drawing code, so no image is delivered and the package declares no asset dependency.  Editorial formatting only, in addition: an AMS \texttt{amsart} preamble that restates the article's own declarations and macros used by the retained lines, namely the environments \texttt{theorem}, \texttt{corollary}, \texttt{lemma} on separate counters, together with the standard packages those lines need.  The retained environments print in this document as Theorem 1, Corollary 1, Lemma 1 in order of appearance, whereas the article prints the same items with the numbering of its own full document.  The complete native source of the article as distributed in the arXiv e-print for this exact version is bundled unchanged as \texttt{sources/208163/source0.tex}.
\begin{equation}\label{Kwong's-condition}
    \operatorname{Ric}_g \geq (n-1)g \quad \text{and} \quad R_g \geq n(n-1)(1+\varepsilon),
\end{equation}

\begin{equation}\label{Bray's-conjecture}
     V_g(M^n) \leq (1 + \varepsilon)^{-\frac{n}{2}} |\mathbb{S}^n|.
\end{equation}

\begin{equation}\label{eq:scaling-mu}
\mu(cg,c\tau)=\mu(g,\tau),
\qquad
\nu(cg)=\nu(g),
\end{equation}

\begin{theorem}\label{thm:MaWang} 
Let $(M^n,g)$ be a connected, closed, smooth Riemannain manifold.
Set
\begin{equation}\label{eq:rho-def}
\rho_g(x):=\max\{0,(n-1)-\lambda_{\min}^g(\Ric_g)(x)\},
\end{equation}
where $\lambda_{\min}^g(\Ric_g)$ denotes  the smallest eigenvalue of the Ricci
tensor.  There exists a constant $\delta_{\MW}(n,p)>0$ such that  if 
\begin{equation}\label{eq:MW-hypothesis}
V_g(M)=|\Sph^n|,
\qquad
\int_M\rho_g^p\,dV_g<\delta_{\MW}(n,p)
\end{equation}
for some positive constant $p>\frac{n}{2}$, 
then the volume-normalized Ricci flow starting at $g$ exists for all forward time
and converges exponentially  to a round metric $g_\infty$
of constant sectional curvature one.
\end{theorem}

\begin{equation}
\rho_g(x):=\max\{0,(n-1)-\lambda_{\min}^g(\Ric_g)(x)\},
\end{equation}

Define \begin{equation}\label{eq:eps-ent}
\eps_{\ent}(n)
=\exp\left(\frac{2\gamma_n}{n}\right)-1
=\frac n{n-1}e^{-1/n}-1.
\end{equation}

\begin{corollary}
\label{cor:entropy-lower}
If $0<\eps\leq\eps_{\ent}(n)$ and \eqref{Kwong's-condition} holds, then
\begin{equation}\label{eq:entropy-sharp-factor}
\nu(g)\geq\nu(g_{\Sph^n})+\log\frac{V_g(M)}{|\Sph^n|}
+\frac n2\log(1+\eps).
\end{equation}
In particular, if
\begin{equation}\label{eq:volume-violation-entropy-0}
V_g(M)>|\Sph^n|(1+\eps)^{-n/2},
\end{equation}
then,
\begin{equation}\label{eq:volume-violation-entropy}
\nu(g)>\nu(g_{\Sph^n}).
\end{equation}
\end{corollary}

\begin{lemma}\label{lem:round-flow-upper}
Suppose $V_g(M)=|\Sph^n|$ and the volume-normalized Ricci flow $g(t)$ starting at
$g$ exists for all $t\geq0$ and converges smoothly  to a metric of constant sectional curvature one. Then
\begin{equation}\label{eq:nu-upper}
\nu(g)\leq\nu(g_{\Sph^n}).
\end{equation}
\end{lemma}

\begin{lemma}\label{lemma:equalrigid}
    Under the assumptions of Lemma \ref{lem:round-flow-upper}, suppose in addition that $\nu(g(0))=\nu(g_{\Sph^n})$, then $$(M, g(0))\cong(\Sph^n, g_{\Sph^n}).$$
\end{lemma}

\begin{theorem}
        Let $(M^n,g)$ be a connected, complete,  smooth Riemannian manifold with $n\geq 3$. There exists a positive constant $\tilde\varepsilon_n>0$ such that, if $0<\varepsilon\leq \tilde\varepsilon_n$ and \eqref{Kwong's-condition} holds,   then \eqref{Bray's-conjecture} holds. {Moreover, the equality in \eqref{Bray's-conjecture} holds if and only if $(M, g)\cong(\Sph^n, (1+\eps)^{-1}g_{\Sph^n})$.}
\end{theorem}

\begin{proof}
First, Bonnet-Myers thereom implies that $M$ is compact. 
Let $\delta_{MW}(n,n)$ be the constant from Theorem \ref{thm:MaWang} and $\varepsilon_{\ent}(n)$ in \eqref{eq:eps-ent}, and define
\begin{equation}\label{eq:epsilon-n-definition}
\tilde\varepsilon_n
:=\min\left\{
\varepsilon_{\ent}(n),
\left(
\frac{\delta_{\MW}(n,n)}{|\Sph^n|(n-1)^n}
\right)^{1/n}
\right\}.
\end{equation}
Both terms inside the minimum are positive, so $\tilde\eps_n>0$.


Fix $0<\eps\leq\tilde \eps_n$ and suppose, for contradiction, that
\begin{equation}\label{eq:contradiction-volume}
V_g(M)>|\Sph^n|(1+\eps)^{-n/2}.
\end{equation}
Since $0<\eps\leq\eps_{\ent}(n)$, Corollary~\ref{cor:entropy-lower} gives
\begin{equation}\label{eq:main-lower}
\nu(g)>\nu(g_{\Sph^n}).
\end{equation}
Bishop's theorem gives $V_g(M)\leq |\Sph^n|$. Define
\begin{equation}\label{eq:main-rescale}
c:=\left(\frac{|\Sph^n|}{V_g(M)}\right)^{2/n},
\qquad h:=cg.
\end{equation}
Thus, one has 
\begin{equation}\label{eq:c-range}
V_h(M)=|\Sph^n|,
\qquad 1\leq c<1+\eps.
\end{equation}
Under constant rescaling, the Ricci tensor viewed as a $(0,2)$ tensor, is unchanged, so
\begin{equation}\label{eq:rescaled-Ricci}
\Ric_h=\Ric_g\geq(n-1)g=\frac{n-1}{c}h.
\end{equation}
The deficit \eqref{eq:rho-def} therefore satisfies
\begin{equation}\label{eq:defect-pointwise}
0\leq\rho_h\leq(n-1)\left(1-\frac1c\right)
<(n-1)\frac{\eps}{1+\eps}.
\end{equation}
Consequently,
\begin{equation}\label{eq:defect-integral}
\int_M\rho_h^n\,dV_h
<|\Sph^n|(n-1)^n\left(\frac{\eps}{1+\eps}\right)^n<\delta_{\MW}(n,n).
\end{equation}
Theorem \ref{thm:MaWang} applies to $h$, and
Lemma \ref{lem:round-flow-upper} gives $\nu(h)\leq\nu(g_{\Sph^n})$. Scale
invariance \eqref{eq:scaling-mu} gives
\begin{equation}\label{eq:main-upper}
\nu(g)=\nu(h)\leq\nu(g_{\Sph^n}),
\end{equation}
contradicting \eqref{eq:main-lower}. Thus \eqref{Bray's-conjecture} holds.

{Now we assume the equality in \eqref{Bray's-conjecture}: 
\begin{equation*}
V_g(M)=(1+\eps)^{-\frac{n}{2}}|\Sph^n|.    
\end{equation*}
The entropy comparison Corollary \ref{cor:entropy-lower} yields
\begin{equation}\label{eq:equalnu}
\nu(g)\geq\nu(g_{\Sph^n})+\log(1+\eps)^{-\frac{n}{2}}
+\frac n2\log(1+\eps)=\nu(g_{\Sph^n}).
\end{equation}
Now by definition in \eqref{eq:main-rescale}, we know $c=1+\eps$ and $h=(1+\eps)g$. In this case,
$$0\leq \rho_h\leq (n-1)\frac{\eps}{1+\eps}.$$
Note the last inequality in \eqref{eq:defect-integral} is strict, we can still apply Theorem \ref{thm:MaWang} to $h$ and finally get \eqref{eq:main-upper}.
Combining \eqref{eq:main-upper} and \eqref{eq:equalnu}, we have $\nu(h)=\nu(g_{\Sph^n}).$ By Lemma \ref{lemma:equalrigid}, we obtain $(M, h)\cong(\Sph^n, g_{\Sph^n})$, and therefor, $(M, g)\cong(\Sph^n, (1+\eps)^{-1}g_{\Sph^n})$. Thus, we complete the proof.}

\end{proof}

\end{document}
